% manuscript_v4_corrected.tex
% TEMPLATE — corrected narrative based on AUC gap decomposition and RAS experiment.
% Key changes from v3:
%   1. Abstract: Attributes AUC gap primarily to cross-cultural factors, not BA.
%   2. New section: AUC gap decomposition with CHARLS raw16 baseline.
%   3. New section: Reduced Aging Score (RAS) experiment.
%   4. Discussion: Revised interpretation — cross-cultural dominates.
%   5. Updated conclusion with evidence hierarchy.

\documentclass[11pt,a4paper]{article}
\usepackage[margin=2.2cm]{geometry}
\usepackage[T1]{fontenc}
\usepackage[utf8]{inputenc}
\usepackage{booktabs}
\usepackage{longtable}
\usepackage{graphicx}
\usepackage{amsmath}
\usepackage{hyperref}
\usepackage{setspace}
\usepackage{microtype}
\usepackage{array}
\usepackage{xcolor}

\setstretch{1.08}
\setlength{\parskip}{4pt}
\setlength{\parindent}{0pt}

\hypersetup{
  colorlinks=true,
  linkcolor=blue,
  urlcolor=blue,
  pdftitle={KOA Risk Prediction: Internal Validation and Cross-Cultural Transportability},
  pdfauthor={BioAge-KOA-Prediction Team}
}

\title{Predicting Symptomatic Knee Osteoarthritis Risk from CHARLS Data:\\
\large Internal Validation, AUC Gap Decomposition, and\\
Cross-Cultural Transportability Assessment on HRS and ELSA}
\author{BioAge-KOA-Prediction Team}
\date{June 2026}

\begin{document}
\maketitle

% ============================================================
\begin{abstract}
% ============================================================
\textbf{Objective.} To develop and internally validate a machine-learning model for
symptomatic knee osteoarthritis (KOA) risk prediction in a CHARLS-derived Chinese cohort,
and to characterise transportability to US (HRS) and UK (ELSA) aging cohorts through a
structured AUC gap decomposition.

\textbf{Methods.} A Random Forest model ($n=300$ trees) was trained on 17
sociodemographic and clinical variables (raw17) from CHARLS Dataset~A
($n=12{,}329$, KOA prevalence $13.3\%$). Internal validation used stratified 5-fold
out-of-fold estimation and an early 80/20 stratified unseen holdout. Transportability was
assessed by applying the model to HRS Wave~13 ($n=20{,}605$, 12 of 17 features) and ELSA
Wave~8 ($n=8{,}416$, 14 of 17 features) using harmonised raw16 feature sets (excluding
Biological Age). To decompose the AUC gap into feature-subset and cross-cultural
components, we evaluated CHARLS on the same raw16 subset (internal baseline) and
additionally constructed a Reduced Aging Score (RAS) from biomarkers common to CHARLS
and ELSA (CRP, HbA1c, SBP, total cholesterol).

\textbf{Results.} The primary model achieved ROC-AUC $0.847$ (internal OOF) / $0.907$
(unseen holdout) with raw17. Restricting CHARLS to raw16 produced ROC-AUC $0.804$ /
$0.869$ — a feature-subset penalty of $\Delta=-0.038$. Applying the same raw16 model to
ELSA and HRS yielded ROC-AUC $0.604$ for both cohorts — a cross-cultural penalty of
$\Delta=-0.265$. The total AUC gap ($-0.303$) decomposes as $13\%$ feature-subset vs.\
$87\%$ cross-cultural. The RAS experiment recovered only $\Delta=+0.006$ AUC,
confirming that biological age proxies contribute minimally to closing the gap.
Isotonic recalibration restored calibration in ELSA (ECE $0.097 \rightarrow 0.005$,
slope $0.271 \rightarrow 0.938$) without improving discrimination, and DCA net benefit
remained positive.

\textbf{Conclusions.} The model generalises poorly across cultures, with outcome proxy
heterogeneity and population differences — not missing Biological Age — as the dominant
degraders. Until a harmonised KOA outcome definition is applied across cohorts and full
raw17 features (including KDM-BA from linked biomarker data) are available, clinical
utility claims remain bounded to the internal CHARLS evidence.

\textbf{Keywords:} symptomatic KOA, machine learning, TRIPOD, calibration, DCA,
transportability, AUC decomposition, HRS, ELSA, cross-cultural replication
\end{abstract}

% ============================================================
\section{Introduction}
% ============================================================
[KEEP existing introduction from v3 — the framing as a TRIPOD-compliant internal
validation study with cross-cultural replication is correct. No changes needed here.]

% ============================================================
\section{Methods}
% ============================================================

[KEEP existing methods sections from v3: Study Objective, Cohort, Feature Set (raw17),
Missing Data, Early Holdout, Primary Model, Class-Imbalance, Evaluation Framework,
Interpretability, Differences vs. Fu et al. (2025). All correct.]

% ADD: AUC Gap Decomposition subsection
\subsection{AUC Gap Decomposition Strategy}
To disentangle feature-subset penalty from cross-cultural attenuation, we conducted
the transportability analysis in three stages:

\begin{enumerate}
  \item \textbf{CHARLS raw16 internal baseline.} The primary model was re-trained and
        evaluated on the CHARLS internal partition using only the 14 features common to
        CHARLS, ELSA, and HRS (raw16: dropping Biological Age, wave, and Time from
        raw17, adding a constant wave identifier). Because both training and evaluation
        occur within CHARLS, this isolates the feature-subset penalty with no
        cross-cultural confound.

  \item \textbf{External application.} The raw16 model trained on CHARLS was applied to
        ELSA Wave~8 and HRS Wave~13 without re-fitting. The gap between CHARLS raw16
        holdout performance and external cohort performance is the cross-cultural
        attenuation, net of the feature-subset penalty.

  \item \textbf{Reduced Aging Score (RAS) sensitivity.} To test whether a harmonised
        biological age proxy could recover the feature-subset gap, we constructed RAS
        from four biomarkers available in both CHARLS and ELSA (hs-CRP, HbA1c, systolic
        BP, total cholesterol). RAS was trained as a linear predictor of CHARLS KDM-BA
        ($R^2=0.417$) and added to the raw16 model for ELSA external validation.
\end{enumerate}

\subsection{Reduced Aging Score (RAS)}
\label{sec:ras}

RAS was constructed as a linear combination of four biomarkers:

\[
\text{RAS} = \beta_0 + \beta_1 \cdot \ln(\text{CRP}) + \beta_2 \cdot \ln(\text{HbA1c})
           + \beta_3 \cdot \ln(\text{TC}) + \beta_4 \cdot \text{SBP}
\]

The coefficients $\beta_{0..4}$ were estimated by ordinary least squares regression of
CHARLS KDM-BA on the four CHARLS biomarker equivalents
(\texttt{crp\_mg.L}, \texttt{Hb A1c}, \texttt{TC\_mg.d L}, \texttt{sbp.mean}).
$R^2$ on CHARLS data was $0.417$ ($n=12{,}329$), indicating that RAS captures
$41.7\%$ of KDM-BA's internal variance.

For ELSA, biomarkers were converted to CHARLS-equivalent units: hs-CRP (mg/L)
log-transformed directly; HbA1c converted from IFCC (mmol/mol) to DCCT (\%) via
$\text{HbA1c}\% = (\text{mmol/mol} / 10.929) + 2.15$, then log-transformed; total
cholesterol converted from mmol/L to mg/dL ($\times 38.67$), then log-transformed;
systolic BP used directly (mmHg). After excluding ELSA respondents with missing
biomarker codes ($-8$, $-2$, $-1$), $2{,}499$ of $8{,}416$ respondents had complete
biomarker data for RAS computation.

% ============================================================
\section{Results}
% ============================================================

[KEEP existing primary model results, reliability diagnostics, DCA, sensitivity branches
from v3 — all still correct.]

% ADD: AUC Gap Decomposition
\subsection{AUC Gap Decomposition: Feature-Subset vs.\ Cross-Cultural}

Table~\ref{tab:auc-decomp} presents the structured decomposition of the
$0.303$ ROC-AUC gap between the CHARLS raw17 holdout and the external cohorts.

\begin{table}[h!]
\centering
\small
\caption{AUC gap decomposition across CHARLS raw17, CHARLS raw16, ELSA, and HRS.
All values are raw (uncalibrated) ROC-AUC except where noted.}
\label{tab:auc-decomp}
\begin{tabular}{lcccccc}
\toprule
Stage & Cohort & Features & $n$ & KOA\% & ROC-AUC & $\Delta$ from baseline \\
\midrule
A  & CHARLS raw17 Holdout  & 17 & 2{,}466  & 13.3 & 0.907 & — (baseline) \\
B  & CHARLS raw16 Holdout  & 15 & 2{,}466  & 13.3 & 0.869 & $-0.038$ (feature-subset) \\
C  & ELSA W8 raw16         & 14 & 8{,}416  & 24.9 & 0.604 & $-0.303$ (total) \\
C' & HRS W13 raw16         & 12 & 20{,}605 & 33.9 & 0.604 & $-0.303$ (total) \\
\cmidrule{6-7}
   & \textbf{Feature-subset component}  & & & & & $\mathbf{-0.038}$ (13\%) \\
   & \textbf{Cross-cultural component}  & & & & & $\mathbf{-0.265}$ (87\%) \\
\bottomrule
\end{tabular}

\vspace{4pt}
\footnotesize{Feature-subset penalty = B $-$ A. Cross-cultural penalty = C $-$ B.
Percentages are of the total A$\rightarrow$C gap ($0.303$).}
\end{table}

The feature-subset penalty (dropping Biological Age, wave, and Time from raw17)
accounts for $13\%$ of the total AUC loss. The remaining $87\%$ is cross-cultural
attenuation occurring between the CHARLS and external cohorts, attributable to
outcome proxy mismatch, population demographic differences, and diagnostic
ascertainment heterogeneity.

Notably, the ELSA and HRS raw16 ROC-AUC values are nearly identical ($0.604$)
despite ELSA having two more harmonised features (Residence and Dyslipidemia),
using a different KOA proxy (ever-diagnosed vs.\ current-wave arthritis), and
drawing from a different population (UK vs.\ US). This convergence supports a
stable transportability ceiling around $\text{AUROC} \approx 0.60$ for raw16
without Biological Age across culturally distinct aging cohorts.

\subsection{Reduced Aging Score (RAS) Experiment}
\label{sec:ras-results}

Table~\ref{tab:ras} presents the effect of adding RAS to the ELSA raw16 model.

\begin{table}[h!]
\centering
\small
\caption{Effect of the Reduced Aging Score on ELSA external validation.
RAS captures 41.7\% of KDM-BA variance in CHARLS.}
\label{tab:ras}
\begin{tabular}{lccc}
\toprule
Model & ROC-AUC & PR-AUC & Brier \\
\midrule
ELSA raw16              & 0.6050 & 0.297 & 0.200 \\
ELSA raw16 + RAS        & 0.6111 & 0.316 & 0.188 \\
\cmidrule{2-4}
$\Delta$ (RAS gain)     & \textbf{+0.0061} & +0.019 & $-0.013$ \\
\bottomrule
\end{tabular}
\end{table}

Adding RAS — the best biological age proxy constructible from publicly available
ELSA biomarkers — recovers only $0.006$ ROC-AUC, or approximately $2\%$ of the
total cross-cultural gap. This negligible gain is consistent with the weak
association between KDM-BA and KOA in CHARLS ($r=0.018$) and the fact that RAS
captures less than half of that already-weak signal.

[KEEP existing HRS and ELSA replication results sections, but update the
interpretation paragraphs to reflect the decomposition above.]

% ============================================================
\section{Discussion}
% ============================================================

\paragraph{Internal validation.}
[KEEP existing internal validation paragraph — still correct.]

\paragraph{AUC gap decomposition and transportability.}
The structured decomposition (Table~\ref{tab:auc-decomp}) reveals that the dominant
driver of performance degradation from CHARLS to ELSA/HRS is \textbf{cross-cultural
shift} ($87\%$), not the exclusion of Biological Age ($13\%$). This finding
contradicts the initial hypothesis — based on SHAP feature importance ranking in the
CHARLS raw17 model — that BA would account for the majority of the external
performance gap. SHAP importance ranks the \textit{within-model} contribution of a
feature to predictions on the training distribution; it does not quantify how much
that feature would contribute to closing a \textit{between-population} performance
gap, which is fundamentally a transportability question.

Three lines of evidence converge on this interpretation:
\begin{enumerate}
  \item \textbf{CHARLS raw16 baseline.} Restricting CHARLS to raw16 (dropping BA,
        wave, and Time) degrades holdout ROC-AUC by only $0.038$ — a modest penalty
        relative to the $0.303$ total gap.
  \item \textbf{ELSA--HRS convergence.} ELSA and HRS produce near-identical ROC-AUC
        ($0.604$) despite different feature availability (14 vs.\ 12 variables),
        KOA proxy definitions (ever-diagnosed vs.\ current-wave arthritis), and
        populations (UK vs.\ US). This convergence is not compatible with a
        BA-centric narrative, which would predict different performance in cohorts
        with different feature completeness.
  \item \textbf{RAS experiment.} Constructing the best possible biological age proxy
        from ELSA biomarkers and adding it to the model recovers only $0.006$
        ROC-AUC — less than one-tenth of the feature-subset penalty and negligible
        relative to the total gap.
\end{enumerate}

The practical implication is that future external validation efforts should
prioritise \textbf{outcome harmonisation} (applying a consistent KOA definition
across cohorts) and \textbf{full raw17 feature availability} (including
KDM-computed Biological Age from linked biomarker panels) over constructing partial
aging proxies from incomplete biomarker sets.

\paragraph{Comparison with the reference study.}
[KEEP existing comparison paragraph — still correct.]

\paragraph{Strengths.}
[KEEP strengths, ADD:] the structured AUC gap decomposition isolating feature-subset
from cross-cultural components; and the RAS sensitivity experiment quantifying the
marginal contribution of harmonised biological age proxies.

% ============================================================
\section{Limitations}
% ============================================================
[KEEP all existing limitations, ADD:]

\begin{itemize}
\item \textbf{Outcome proxy heterogeneity is the primary limitation.} The AUC gap
      decomposition shows that $87\%$ of performance loss is cross-cultural, with
      outcome definition mismatch (CHARLS knee-localised pain vs.\ HRS/ELSA stooping
      difficulty) likely the largest single contributor. Until a harmonised KOA
      definition is applied across all three cohorts, the true transportability of
      the raw17 model cannot be estimated.
\item \textbf{RAS is an imperfect BA proxy.} The Reduced Aging Score captures only
      $42\%$ of KDM-BA variance and is based on 4 of the 8 KDM biomarkers. It should
      not be interpreted as a substitute for full KDM-BA. The RAS experiment serves
      as a sensitivity analysis demonstrating the ceiling of what is achievable with
      publicly available ELSA data.
\end{itemize}

% ============================================================
\section{Conclusion}
% ============================================================

This study demonstrates a methodologically rigorous internal-validation pipeline for
symptomatic KOA prediction in a large Chinese nationally-representative cohort. The
pre-specified primary model achieves strong discrimination (ROC-AUC $0.907$ on the
unseen holdout), acceptable calibration, and positive DCA net benefit in the
clinically relevant $10\%$--$30\%$ risk threshold band.

A structured AUC gap decomposition reveals that cross-cultural transportability
losses ($87\%$ of the total gap) dominate over feature-subset penalties ($13\%$).
The near-identical performance of the raw16 model on ELSA and HRS (ROC-AUC
$\approx 0.60$) — despite different feature availability, KOA proxy definitions, and
populations — establishes a conservative transportability ceiling. A Reduced Aging
Score constructed from biomarkers common to CHARLS and ELSA recovers negligible AUC
($+0.006$), confirming that partial biological age proxies are insufficient to bridge
the cross-cultural gap.

These findings reframe the external validation challenge: the primary obstacle is not
missing Biological Age but rather outcome definition heterogeneity and cross-cultural
population differences. Future work should prioritise (i) harmonised KOA outcome
definitions across cohorts, (ii) full raw17 feature availability including
KDM-computed BA from linked biomarker panels, and (iii) prospective validation in
cohorts where the same KOA definition can be applied as in CHARLS.

% ============================================================
\section*{Data Availability}
% ============================================================
[KEEP existing data availability section — still correct.]

% ============================================================
\begin{thebibliography}{99}
% ============================================================
[KEEP all existing references, ADD if needed for RAS/AUC decomposition citations.]
\end{thebibliography}

\end{document}
